\begin{frame}
	\frametitle{Linha do tempo: da quantização discreta às LLMs de 1 bit}
	\centering
	\resizebox{.98\linewidth}{!}{%
	\begin{tikzpicture}[
		yr/.style={draw=bitnetColor, fill=bitnetColor!15, rounded corners=2pt, minimum height=0.9cm, minimum width=1.6cm, font=\tiny\sffamily, align=center},
		lb/.style={font=\tiny\sffamily, align=center, text width=1.7cm}
		]
		\draw[-{Latex[length=2mm]}, line width=0.6mm, draw=neutralColor] (-0.3,0) -- (12.3,0);

		\node[yr] (n1) at (0,0) {2015};
		\node[lb, below=2pt of n1] {BinaryConnect \cite{courbariaux2016binarized}};

		\node[yr] (n2) at (2.4,0) {2016};
		\node[lb, above=2pt of n2] {BNN \\ XNOR-Net \cite{rastegari2016xnor} \\ TWN \cite{li2016ternary}};

		\node[yr] (n3) at (4.8,0) {2017};
		\node[lb, below=2pt of n3] {TTQ \cite{zhu2017trained}};

		\node[yr] (n4) at (7.2,0) {2022};
		\node[lb, above=2pt of n4] {LLM.int8() \cite{dettmers2022llm}};

		\node[yr] (n5) at (9.6,0) {2023};
		\node[lb, below=2pt of n5] {GPTQ \cite{frantar2023gptq} \\ SmoothQuant \cite{xiao2023smoothquant} \\ BitNet \cite{wang2023bitnet}};

		\node[yr, fill=bitnetColor!45] (n6) at (12,0) {2024};
		\node[lb, above=2pt of n6] {\textbf{BitNet b1.58} \cite{ma2024era}};
	\end{tikzpicture}%
	}
	\vspace{2pt}
	\begin{itemize}\small
		\item 2015--2017: quantização extrema nasce em \textbf{visão computacional} (redes pequenas, treinadas do zero).
		\item 2022--2023: quantização \textbf{pós-treino} (PTQ) vira estratégia dominante para LLMs já treinados.
		\item 2023--2024: BitNet retoma a ideia de \textbf{treinar do zero} em baixa precisão --- agora em escala de bilhões de parâmetros.
	\end{itemize}
\end{frame}
